\section{Scaling}
\label{sec:scaling}

\begin{question}
  \textbf{How does throughput per GPU change from one switch up to 264 GPUs, and where does efficiency
  drop?} The work per GPU is fixed (weak scaling), so perfect scaling is a flat line.
\end{question}

\Cref{fig:scaling} shows tokens/s per GPU against the number of GPUs, one line per cluster, with the
change from the smallest to the largest run written at the end of each line. There is no ``efficiency''
axis: the only runs without inter-node traffic are single-node runs, which use a 16-layer model, so no
full-model run can serve as a 100\,\% reference. The dashed lines are a \emph{no-network ceiling}
instead: the single-node MFU converted to full-model tokens/s. It is an estimate: it assumes MFU carries
over from 16 to 64 layers, and the 16-layer model also gathers 4$\times$ fewer parameters per step.

\begin{figure}[h]
  \centering\fig{weak_scaling}
  \caption{Weak scaling of Qwen3.8-27B ($3\times4096$ tokens per GPU, FP32 moments). One point per run
    (all placements), lines through the medians. Dashed: no inter-node traffic (est.).
    \syntheticnote{weak_scaling}}
  \label{fig:scaling}
\end{figure}

\placeholder{throughput at 12, 24, 48 GPUs (measured) and at full scale; where the curve bends.}

\paragraph{More nodes vs.\ crossing switches.} Throughput per GPU drops as a job grows (\cref{fig:scaling}).
\Cref{fig:placement-scaling} separates two causes up to 12 nodes, the largest job one Rorqual switch holds.
The one-switch line (S) is the cost of more nodes alone, with no traffic through the spine; its band is the
run-to-run spread. The spread placements on the same axis, 3 per switch (H) and 1 per switch (O), show what
crossing switches adds. At 3 nodes, 3 per switch is a single 3-node leaf, so it is not shown as spreading.
Instead, a leaf-to-leaf control (two one-switch runs on different leaves) shows how much leaves differ by
themselves. All clusters share one axis: each is relative to its own one-switch run at 3 nodes. 1 per
switch is limited by the number of leaves (16 on Rorqual, 8 GPU leaves on Fir, 5 on Tamia). Beyond 12
nodes, the two causes cannot be separated on Rorqual (no one-switch reference); the allocation choice at
that scale is in \cref{fig:placement-at-scale}.

\begin{figure}[h]
  \centering\fig{placement_scaling}
  \caption{Throughput per GPU by number of nodes for each placement (marker) and cluster (color), relative to
    the cluster's one-switch run at 3 nodes; one campaign per cluster. Band: run-to-run spread around one
    switch. Shaded (Rorqual): what crossing switches adds. Gray X: leaf-to-leaf control.
    \syntheticnote{placement_scaling}}
  \label{fig:placement-scaling}
\end{figure}

\placeholder{share of the drop at 12 nodes due to node count vs.\ crossing switches, per cluster.}

\paragraph{Where the step time goes.} \Cref{fig:decomposition} splits the step time into compute and exposed
communication. The compute part is the single-node step time at the same per-GPU batch, scaled by the
ratio of FLOPs per token between the full and the 16-layer model. It is an estimate, and so is the
communication part derived from it.

\begin{figure}[h]
  \centering\fig[0.85\linewidth]{step_decomposition}
  \caption{Step time decomposition by cluster and scale. Hatched: compute (est.); shaded: exposed
    communication (est.), the remainder of the measured step time. \syntheticnote{step_decomposition}}
  \label{fig:decomposition}
\end{figure}

\placeholder{share of exposed communication on Rorqual vs.\ Fir at equal scale; a measured split would
need gradient-accumulation runs at multi-node scale ($t_\text{step} = k\,t_\text{micro} +
t_\text{overhead}$).}
